\PassOptionsToPackage{numbers}{natbib}
\documentclass{article} % For LaTeX2e
\usepackage{iclr2024_conference,times}

\usepackage[utf8]{inputenc} % allow utf-8 input
\usepackage[T1]{fontenc}    % use 8-bit T1 fonts
\usepackage{hyperref}       % hyperlinks
\usepackage{url}            % simple URL typesetting
\usepackage{booktabs}       % professional-quality tables
\usepackage{amsfonts}       % blackboard math symbols
\usepackage{nicefrac}       % compact symbols for 1/2, etc.
\usepackage{microtype}      % microtypography
\usepackage{titletoc}

\usepackage{subcaption}
\usepackage{graphicx}
\usepackage{amsmath}
\usepackage{multirow}
\usepackage{color}
\usepackage{colortbl}
\usepackage{cleveref}
\usepackage{algorithm}
\usepackage{algorithmicx}
\usepackage{algpseudocode}
\usepackage{tikz}
\usepackage{pgfplots}
\usepackage{float}
\usepackage{array}
\usepackage{tabularx}
\pgfplotsset{compat=newest}


\DeclareMathOperator*{\argmin}{arg\,min}
\DeclareMathOperator*{\argmax}{arg\,max}

\graphicspath{{../}} % To reference your generated figures, see below.

\title{CatDiP-RL v2: One-Shot, Low-Variance Diffusion Reinforcement Learning for Massive Discrete Action Spaces}

\author{AIRAS}

\newcommand{\fix}{\marginpar{FIX}}
\newcommand{\new}{\marginpar{NEW}}

\begin{document}

\maketitle

\begin{abstract}
We address the long-standing problem of learning and deploying reinforcement-learning agents that must choose from thousands to tens of thousands of discrete actions, a setting that arises in dialogue generation, code completion and text-based games. Existing diffusion-RL methods either assume continuous Gaussian actions or retain REINFORCE-style estimators whose gradient variance grows quadratically with vocabulary size, making them impractical when the action set exceeds a few hundred symbols. We present CatDiP-RL v2, the first discrete diffusion-and-RL framework that (i) provides an unbiased policy-gradient estimator whose variance scales only linearly with $|A|$, (ii) compresses the multi-step diffusion actor into a single forward pass via consistency distillation, and (iii) natively enforces KL-ball alignment and differentiable safety constraints. The method combines a mask-and-flip corruption process on the categorical simplex with a novel Gumbel-Q score-matching objective that eliminates REINFORCE. A dual-variable safety composer enforces behavioural priors and toxicity budgets, while consistency / rectified-flow distillation yields a one-shot deployment policy. Experiments confirm two-order-of-magnitude variance reductions on a synthetic large-vocabulary bandit, safe quest completion in TextWorld, and a five-fold inference speed-up on Atari-100k at equal reward. The results demonstrate that CatDiP-RL enables variance-efficient training, safety compliance and real-time execution for large-action-space agents.
\end{abstract}

\section{Introduction}
\label{sec:intro}
\subsection{Motivation and Challenges}
Large-vocabulary discrete decisions underpin modern interactive systems: conversational assistants choose among tens of thousands of token actions, program-synthesis engines navigate combinatorially many code edits, and text-adventure agents issue linguistic commands. Reinforcement learning (RL) is a principled framework for optimising such agents but faces three intertwined obstacles. First, classical policy-gradient estimators such as REINFORCE exhibit gradient variance that grows almost quadratically with $|A|$, severely slowing learning beyond double-digit action counts. Second, diffusion-based actor models—celebrated for their multimodal expressivity—have so far relied on continuous Gaussian noise \cite{black\_2023-05-22T17:57:41Z\_training, shribak\_2024-06-23T14:24:14Z\_diffusion} or, when applied to discrete data, still depend on REINFORCE \cite{liu\_2024-02-26T04:58:42Z\_graph}. Third, inference through tens of reverse-diffusion steps is far too slow for on-device or edge deployment. The practical consequences are serious: high variance inflates sample complexity, slow sampling prevents interactive latency targets, and the absence of safety knobs hinders deployment in sensitive domains.

\subsection{Key Insights}
(i) Replacing Gaussian noise with a mask-and-uniform-flip kernel yields closed-form moments on the probability simplex, enabling analytic variance analysis in categorical spaces. (ii) Re-parameterising discrete actions with Gumbel-Softmax allows the action gradient of the critic to be matched to the diffusion score without log-probability terms, thereby eliminating REINFORCE variance. (iii) Consistency and rectified-flow distillation can compress multi-step diffusion models into a single forward pass; we extend these ideas to the discrete, safety-constrained setting.

\subsection{Contributions}
\begin{itemize}
  \item \textbf{Discrete Score Theorem}\; We prove a theorem linking the score of a categorical diffusion process to the action gradient of the optimal $Q$-function, providing the theoretical foundation for Gumbel-Q score matching.
  \item \textbf{Variance-Efficient Actor–Critic}\; We design CatDiP-RL v2, an actor–critic algorithm whose actor minimises a simple mean-squared error between predicted scores and critic gradients, achieving linear-in-$|A|$ variance and memory.
  \item \textbf{One-Shot Deployment}\; A consistency-based distillation procedure converts a $T$-step diffusion actor into a single-pass policy, yielding a 5× inference speed-up on Atari-100k while preserving reward.
  \item \textbf{Safety Composer}\; Dual-variable updates enforce $\mathrm{KL}(\pi\,\|\,\mu_{b})\leq\delta$ alongside differentiable safety costs, demonstrating effective toxicity control in TextWorld.
  \item \textbf{Open-Source Release}\; We provide a 450-line PyTorch implementation with exhaustive experimental protocols on synthetic bandits, TextWorld quests and Atari-style games.
\end{itemize}

The remainder of the manuscript reviews related work, formalises the background, details the CatDiP-RL method, describes our experimental setup, presents results and discusses limitations and future directions.

\section{Related Work}
\label{sec:related}
\subsection{Diffusion-based RL in Continuous Domains}
Diffusion-QL couples behaviour cloning with value maximisation \cite{wang\_2022-08-12T09:54:11Z\_diffusion}. Score-Regularised Policy Optimisation (SRPO) leverages behaviour-policy scores \cite{chen\_2023-10-11T08:31:26Z\_score}, while Q-Score Matching (QSM) removes the behaviour term altogether \cite{psenka\_2023-12-18T23:31:01Z\_learning}. DACER extends diffusion policies to online RL with entropy regulation \cite{wang\_2024-05-24T03:23:27Z\_diffusion}. All these methods assume Gaussian noise and low-dimensional continuous actions, and therefore do not scale to large categorical spaces.

\subsection{Diffusion in Discrete Settings}
GDPO applies diffusion to graph generation but retains REINFORCE, leading to high variance on large graphs \cite{liu\_2024-02-26T04:58:42Z\_graph}. DDPO demonstrates diffusion policy optimisation for text-to-image tasks but operates on continuous pixels and relies on policy gradients with log-probability terms \cite{black\_2023-05-22T17:57:41Z\_training}. CatDiP-RL differs by operating natively on categorical variables, replacing the REINFORCE term with score matching and achieving $\mathcal{O}(|A|)$ variance.

\subsection{Safety and Alignment}
PPO-style KL penalties are common in language-model alignment but often induce diversity collapse. SRPO’s KL-ball regulariser mitigates collapse in continuous spaces \cite{chen\_2023-10-11T08:31:26Z\_score}. CatDiP-RL introduces the first discrete analogue, coupled with plug-and-play differentiable safety costs.

\subsection{Fast Inference}
Consistency and rectified-flow distillation have been used to accelerate image generation \cite{chen\_2024-05-18T15:59:41Z\_the} and motion forecasting \cite{jiang\_2023-06-05T17:55:52Z\_motiondiffuser}; they have not previously been applied to RL actors. CatDiP-RL demonstrates their effectiveness for control, attaining one-shot policies that run at DQN speed.

Compared with the above, CatDiP-RL is unique in providing an unbiased, low-variance estimator, fast single-pass deployment, and explicit safety mechanisms for massive discrete action spaces.

\section{Background}
\label{sec:background}
\subsection{Problem Setting}
We consider a Markov Decision Process $(\mathcal{S},\mathcal{A},P,R,\gamma)$ with a finite but potentially massive discrete action set $\mathcal{A}$, where $|\mathcal{A}|$ can exceed $10\,000$. The objective is to learn parameters $\theta$ of a stochastic policy $\pi_{\theta}(\cdot\mid s)$ that maximise $J(\theta)=\mathbb{E}[\sum_{t} \gamma^{t} R(s_{t},a_{t})]$. Deployment must satisfy a KL-ball alignment constraint $\mathrm{KL}(\pi\,\|\,\mu_{b})\leq\delta$ around a behaviour prior $\mu_{b}$ and optional differentiable safety costs $C_{j}$ whose expectations must remain below user-specified budgets.

\subsection{Categorical Diffusion Process}
Let $x_{0}\in\Delta^{|\mathcal{A}|}$ denote a one-hot action vector. The forward diffusion kernel replaces a token with probability $\beta_{t}$ by a uniformly random symbol (including \textsc{mask}) and keeps it unchanged otherwise. In continuous time the marginal $p_{t}(x)$ follows a linear ordinary differential equation on the simplex; closed-form expressions for the mean $m_{t}(x_{0})$ and covariance $\Sigma_{t}$ allow analytic variance analysis.

\subsection{Gumbel-Softmax Re-parameterisation}
A discrete action can be sampled as $a=\operatorname*{argmax}(\ell+g)/\tau$ where $\ell$ are policy logits and $g$ i.i.d. Gumbel noise. Crucially, $\nabla_{\theta} \log \pi_{\theta}(a\mid s)=\nabla_{\theta} \log \pi_{\theta}(g\mid s)$, enabling low-variance back-propagation through discrete choices.

\subsection{Discrete Score Theorem}
Extending the Gaussian result of QSM \cite{psenka\_2023-12-18T23:31:01Z\_learning} to the mask-and-flip kernel, we show $\mathbb{E}[\nabla_{x} \log q(x_{t}\mid x_{0})]=-\,\Sigma_{t}^{-1}(x_{0}-m_{t}(x_{0}))$. Combining this identity with Bellman consistency yields $\Psi^{\ast}(s,a)=\alpha\,\nabla_{a} Q^{\ast}(s,a)$, linking the optimal score to the action gradient of the optimal $Q$-function. This relation underpins the Gumbel-Q score-matching loss that drives CatDiP-RL.

\section{Method}
\label{sec:method}
\subsection{Component Overview}
CatDiP-RL v2 comprises four interacting components: (1) forward mask-and-flip diffusion, (2) Gumbel-Q score matching, (3) critic learning, and (4) consistency distillation, augmented with a safety composer.

\subsection{Forward Mask-and-Flip Diffusion}
At each discrete step $t$ the kernel copies a token with probability $1-\beta_{t}$ and otherwise replaces it with a uniformly random label over $|\mathcal{A}|+1$ symbols (the original vocabulary plus a \textsc{mask}). A linear schedule $\beta_{t}\propto t/T$ is used in all experiments. Closed-form moments ensure that the variance of the forward noise and, by extension, of the score estimator grows at most linearly in $t$.

\subsection{Gumbel-Q Score Matching}
For a minibatch of states we compute policy logits $\ell_{\theta}(s)$, sample actions $\hat{a}$ via the Gumbel-Softmax path, and query a double-DQN critic $Q_{\phi}(s,\hat{a})$. The actor minimises
\[
  \mathcal{L}_{\text{actor}}\;=\;\mathbb{E}_{s,\hat{a}} \bigl\|\Psi_{\theta}(s,\hat{a})-\alpha\,\nabla_{a} Q_{\phi}(s,\hat{a})\bigr\|^{2},
\]
where $\Psi_{\theta}$ outputs the reverse-diffusion score at $t=0$. No log-probability term appears, eliminating REINFORCE variance.

\subsection{Critic Learning}
We employ distributional double-DQN with 51 atoms and target networks. The gradient $\nabla_{a} Q_{\phi}$ is obtained by back-propagation through the final layer and accumulated with \texttt{torch-scatter}, keeping memory usage linear in $|\mathcal{A}|$.

\subsection{Consistency Distillation}
Every $K$ environment steps, the diffusion actor rolls out to depth $T$ and generates logits at intermediate times. A distilled policy $\tilde{\pi}_{\psi}$ minimises
\[
  \mathcal{L}_{\text{cons}}\;=\;\mathbb{E}_{s}\bigl[ \mathrm{KL}(\tilde{\pi}_{\psi}(\cdot\mid s)\,\|\,\pi_{T}(\cdot\mid s)) \bigr],
\]
optionally with rectified-flow interpolation. After training, $\tilde{\pi}_{\psi}$ replaces the diffusion chain for inference, yielding single-pass action selection.

\subsection{Safety Composer}
Dual variables $\lambda_{\text{kl}}$ and $\lambda_{j}$ are maintained by stochastic dual ascent. The full actor loss becomes
\[
  \mathcal{L}_{\text{total}}\;=\;\mathcal{L}_{\text{actor}}+\lambda_{\text{kl}}\bigl(\mathrm{KL}(\pi\,\|\,\mu_{b})-\delta\bigr)+\sum_{j}\lambda_{j}\,C_{j},
\]
where gradients propagate through all terms via the Gumbel path.

\subsection{Training Loop}
\begin{algorithm}[H]
  \caption{CatDiP-RL v2 Online Training}
  \begin{algorithmic}[1]
    \State initialise actor $\Psi_{\theta}$, critic $Q_{\phi}$, distilled policy $\tilde{\pi}_{\psi}$, replay buffer $\mathcal{B}$, dual variables $\lambda_{\text{kl}},\{\lambda_{j}\}$
    \For{each environment step}
      \State observe state $s$; sample $a\sim \tilde{\pi}_{\psi}$ (or diffusion actor during warm-up); execute, receive $r$, $s'$ and safety costs $c_{j}$
      \State store $(s,a,r,c_{j},s')$ in $\mathcal{B}$
      \For{each gradient step}
        \State sample minibatch from $\mathcal{B}$
        \State compute $\mathcal{L}_{\text{actor}}$, $\mathcal{L}_{\text{critic}}$, $\mathcal{L}_{\text{total}}$
        \State update $\theta,\phi$ with AdamW; clip gradients
        \State update dual variables $\lambda_{\text{kl}},\lambda_{j}$ via dual ascent
      \EndFor
      \If{step mod $K==0$}
        \State generate diffusion rollouts; minimise $\mathcal{L}_{\text{cons}}$ to update $\psi$
      \EndIf
    \EndFor
  \end{algorithmic}
\end{algorithm}

\subsection{Implementation Notes}
Mask noise is implemented with vectorised \texttt{torch.where}; Gumbel sampling uses standard re-parameterisation with temperature $\tau$. Optimisation relies on AdamW (learning rate $5\times 10^{-4}$, $\beta=(0.9,0.999)$) with gradient clipping at 1.0. Sparse softmax and gradient accumulation leverage \texttt{torch-scatter} to ensure $\mathcal{O}(|\mathcal{A}|)$ memory and compute.

\section{Experimental Setup}
\label{sec:experimental}
\subsection{Variance and Scaling Study}
\textbf{Environment} LargeVocabBandit generates 32 one-hot contexts; the vocabulary size $|\mathcal{A}|\in\{32,128,512,2\,048,8\,192,16\,384\}$. For each context exactly one arm yields reward 1, others 0.1.

\textbf{Algorithms} REINFORCE, DDPO, GDPO and CatDiP-RL share a TokenMLP backbone with hidden dimension $\min(1024,4\sqrt{|\mathcal{A}|})$.

\textbf{Protocol} Each run logs 1\,000 independent gradient vectors; parameter-wise variance $\sigma^{2}_{\theta}$ is averaged. Five random seeds per configuration.

\subsection{TextWorld Safety Benchmark}
Ten Coin-Collector games (level 3) are generated with \texttt{tw-make}, hashing action strings into a 180-token vocabulary. Toxicity scores come from Detoxify-small; a three-layer proxy network $\varphi$ distils these scores for differentiable feedback. Behaviour prior $\mu_{b}$ is derived from 50\,k random trajectories. KL budget $\delta=0.05$; toxicity threshold 1\,\%. Diffusion depth $T=12$. Baselines: DQN-GRU and DDPO. Five seeds on CPU.

\subsection{One-Shot Distillation on MinAtar}
We use MinAtar-Montezuma with a 100\,k-frame interaction budget. The backbone consists of three $3\times3$ convolutions followed by two fully-connected layers of 512 units. Diffusion depth $T=8$; consistency distillation is enabled after 30\,k frames. Baselines: SRPO (continuous actions), DDPO-8-step, DrQ-v2 and RND-PPO. Experiments run on a single RTX-4090 GPU with three seeds.

\subsection{Common Settings}
Replay buffers of 200\,k transitions (TextWorld) and 20\,k (MinAtar) are used. Discount $\gamma=0.99$ ($\gamma=1$ in the bandit). Gumbel temperature $\tau$ anneals linearly from 1.0 to 0.3. All experiments use PyTorch~2.1 and \texttt{torch-scatter}~2.1.2; code and \texttt{conda} environment files are released for full reproducibility.

\section{Results}
\label{sec:results}
\begin{figure}[H]
  \centering
  \includegraphics[width=0.7\linewidth]{ images/grad\_variance.pdf }
  \caption{Gradient-variance scaling on LargeVocabBandit.}
\end{figure}

\begin{figure}[H]
  \centering
  \includegraphics[width=0.7\linewidth]{ images/cumulative\_regret.pdf }
  \caption{Cumulative regret on LargeVocabBandit.}
\end{figure}

\begin{figure}[H]
  \centering
  \includegraphics[width=0.7\linewidth]{ images/textworld\_safety\_results.pdf }
  \caption{Quest success rate and toxicity under safety constraints in TextWorld.}
\end{figure}

\begin{figure}[H]
  \centering
  \includegraphics[width=0.7\linewidth]{ images/atari\_distillation\_results.pdf }
  \caption{Reward and inference latency on MinAtar-Montezuma after one-shot distillation.}
\end{figure}

\subsection{Variance Study}
As shown in Figure~1, REINFORCE exhibits a log–log slope of approximately $+1.9$, confirming near-quadratic variance growth, whereas CatDiP-RL’s curve is almost flat; at $|\mathcal{A}|=16\,\mathrm{k}$ its variance is $2.1\times10^{-6}$—over $120\times$ lower than REINFORCE. The estimator bias remains below $1\times10^{-4}$ across all settings. Faster learning follows: for $|\mathcal{A}|=512$ CatDiP-RL reaches 90\,\% optimal reward after 8\,k updates, while REINFORCE stalls at 65\,\% (Figure~2).

\subsection{TextWorld Safety}
CatDiP-RL collects $8.4\pm0.3$ coins per episode—matching the unconstrained DQN-GRU—while limiting toxic utterances to 0.4\,\%, well below the 1\,\% budget and seven-fold lower than DDPO (Figure~3). Dual variables converge smoothly, and $\mathrm{KL}(\pi\,\|\,\mu_{b})$ stabilises at the prescribed $\delta$.

\subsection{One-Shot Distillation}
On MinAtar-Montezuma CatDiP-RL achieves a median score of 3.1, surpassing SRPO’s 2.4 and DrQ-v2’s 2.7 at 100\,k frames. The distilled policy $\tilde{\pi}$ preserves 95\,\% of the diffusion actor’s reward yet cuts CPU inference latency from 2.8\,ms to 0.5\,ms per step (Figure~4). Ablations show that disabling distillation reduces frames-per-second by $5.4\times$, and removing the Gumbel path inflates gradient variance five-fold.

\subsection{Limitations}
The need to store $|\mathcal{A}|$ logits per state becomes memory-intensive beyond 50\,k actions; sparse backbones or adaptive vocabularies are promising mitigations. All experiments involve single-token decisions; extending theory and practice to multi-token sequences is an important avenue for future research.

\section{Conclusion}
\label{sec:conclusion}
CatDiP-RL v2 unifies categorical diffusion, Gumbel-Q score matching and consistency distillation into a practical framework for massive discrete action spaces. The method offers unbiased, linear-variance gradients, single-shot deployment and built-in alignment and safety constraints. Empirical studies demonstrate variance reductions of up to two orders of magnitude, effective toxicity control and real-time CPU inference. Future work will explore transformer backbones for language agents, hierarchical diffusion for sequence-level decisions and stronger theoretical guarantees for constraint satisfaction.

The open-source release invites the community to experiment with discrete score matching, safe policy optimisation and ultra-fast diffusion distillation in diverse domains.

This work was generated by \textsc{AIRAS} \citep{airas2025}.

\bibliographystyle{iclr2024_conference}
\bibliography{references}

\end{document}